% Editable transcription; verify reused text against the original PDF.
\section{Individual differences in intelligence}

One question posed at the start of the paper, that has not yet been
addressed, concerns the issue of individual differences.
The intelligence researcher traditionally focuses on individual
differences in cognitive functioning, and not so much on the cognitive
architecture of the mind.
Is AI then relevant for individual differences research?
We can now answer this question with a definitive yes.

The first somewhat trivial reply is that a deeper understanding of
how the human mind operates informs us about possible sources of
individual differences in intelligence, but modern AI offers more
interesting and explicit insights into individual differences in
(artificial) intelligence.

Chess, and in particular the previously discussed Leela Zero Chess
models, again provides an interesting example.
What has to be done before Leela Chess Zero can play superhuman chess?
First, a network architecture, defined by the number of layers and
filters of the neural network, is chosen.
The choice depends on the hardware and computation time available,
where the largest networks are trained for months on an extensive
and powerful array of computer hardware.
Second, choices about the training methods are made.
Official variants use reinforcement learning and self-play and learn
largely from the ground up, but other variants incorporate supervised
learning and rich knowledge bases (e.g. on openings) to improve
performance.
According to the Leela chess website the LCZero team has produced
hundreds of thousands of AI chess playing systems.

One could thus say that the development of all of these Leela chess
variants mirrors the case of people in intelligence research.
Humans and Leela chess both have information processing systems
of very high complexity, to such a degree that we do not really
understand their mechanisms.
We could say that both nature (the network architecture and the
hardware used to train the system) and nurture (training regime)
are involved.
And we have individual differences: Which variants of Leela chess
are actually the best?

The strength of chess engines is measured by the Elo rating system,
which is mathematically almost identical to the Rasch model, the
root model of item response theory that forms the basis of the best
measurement models in intelligence research.
The Elo rating system suffices when deciding which chess system
is best.
However, Elo ratings provide a one-dimensional measurement of
strength, while the many Leela variants probably differ in other
dimensions too.
Some may excel in speed chess, some may be exceptionally good
at end games or tactical positions.
Some Leelas might be quite adventurous and others rather conservative.
Insight into these individual differences may be best achieved
by devising tests with subtests consisting of items meant to measure
these sub-abilities in chess, akin to what has been done in human
chess research (van der Maas \& Wagenmakers, 2005).

The factor analytic study of these artificial individual differences
is of interest too, not only to understand Leela zero chess better,
but to also see what the factor analytic approach actually brings
us in such a case.
Suppose, and we think this is quite probable, we again find a general
factor, what then does this mean?
Is it a formative or a reflective factor
(van der Maas, Kan, Marsman, \& Stevenson, 2017)?
Does this g-factor represent one underlying source of chess strength
(computing speed, for instance), or does it summarize the collective
force of the many building blocks of Leela?
Note that a part of the success of current AI lies in the combination
of techniques such as deep learning, reinforcement learning and
powerful hardware.
Having such a mirror system creates new possibilities to study
intelligence and also raises interesting questions regarding
individual differences research.

Similarly, as modern AI systems are systems that learn, it is
interesting to rethink questions regarding the development
of intelligence.
One example might be the study of the third source hypothesis
(Kan, Ploeger, Raijmakers,
